\begin{center}
\LARGE
Formal Lie Series and Applications to the Hamiltonian
                           Mechanics

\end{center}

\begin{center}
\Large
Nikolay N.Vasiliev

\end{center}

\noindent
Date:  July 19th (Friday)\\
Time:  17:00-17:20\\

\begin{center}
\large
Abstract
\end{center}

We consider different applications of formal Lie series to the hamiltonian mechanics. First
example is the symplectic integration. Construction of most known symplectic methods, like
Yoshida method is based on iterating process of superpositions of the exponential maps which can
be represented as Lie series. Usually such superpositions can be evaluated on the base of BCH
formula. We show how to construct implicit symplectic methods for numerical integration of
hamiltonian ODE if the hamiltonian does not have the separated form $H(P,Q)=U(P)+V(Q)$.
Another problem which is also connected with application of symplectic methods is calculating of
perturbation of hamiltonian due to concrete symplectic method for numerical integration. Some
results about modeling of Henon-Heyles system via symplectic and nonsymplectic methods are
also presented. We show how to use the procedure of normalisation based on the Lie series in order
to construct the formal integral for the symplectic approximations of the Henon-Heiles system.
We also discribe some algorithms for processing with formal Lie series, the problem of computer
constructing of special basses in free Lie algebra and methods for representation of Lie series in
the computer memory which allow effectively realize such algorithms. 

